\paragraph{Task Description.} We apply \ouralgo to the ALE-Bench LITE \citep{imajuku2025ale} benchmark, a collection of 10 competitive programming contests hosted by AtCoder and designed to test the performance of LLMs on heuristic problems. Here, we explore whether \ouralgo can successfully improve high-performing solutions discovered by LLMs. We leverage the best programming solution discovered by ALE-Agent \citep{imajuku2025ale} as an initial program for each problem and apply \ouralgo to improve on top of it. We run \ouralgo for 50 generations, leveraging the score calculated on the public test set as the fitness function. Afterwards, we submit the best solution to the private test set and report the score.

\paragraph{\ouralgo's Discovery Dynamics.} \ouralgo is able to improve the solutions discovered by ALE-Agent by approximately 2.3\% across the 10 tasks on average (\Cref{fig:results_alebench}). Furthermore, on one task, \texttt{ahc039}, the combination of \ouralgo with ALE-Agent resulted in the second place submission on the \href{https://atcoder.jp/contests/ahc039/submissions/68865341}{AtCoder leaderboard} if they had participated. While these improvements resulted from detailed implementation improvements, we observe that the proposed changes by \ouralgo remained algorithmically close to the original ALE-Agent's initialization solution.

\begin{figure}[h!]
    \centering
    \includegraphics[width=0.975\textwidth]{figures/f6_results_alebench.pdf}
    \caption{\textbf{\ouralgo for Improving ALE-Bench solutions.} \textit{Left:} \ouralgo improves the solutions discovered by ALE-Agent by $\sim2.3\%$. \textit{Right:} On one task, \texttt{ahc039}, the solution improved from 5th to 2nd place submission on the AtCoder leaderboard if it had participated in the contest.}
    \label{fig:results_alebench}
    \vspace{-15pt}
\end{figure}

\paragraph{\ouralgo's Discovered Solution.}
We focus on two tasks to illustrate the discovered improvements of \ouralgo, \texttt{ahc039} and \texttt{ahc025}.
The objective of \texttt{ahc039} is to find an optimal, axis-aligned polygon to maximize the number of mackerels it contains minus the number of sardines, subject to given constraints. The base solution by ALE-Agent applies simulated annealing with kd-tree data structure (5th, 2880 performance). \ouralgo further improved the solution (2nd, 3140 performance) by introducing modifications such as caching the validation process and enhancing neighborhood operators. For the caching, the kd-tree was augmented to cache subtree statistics, including bounding boxes and fish counts, at each node. For the neighborhood operators, a novel ``targeted edge move'' was introduced, which heuristically identifies a misclassified fish (e.g., a mackerel outside the polygon) and greedily moves the nearest edge to correct its state. These changes strengthened the directionality of the search.
For \texttt{ahc025}, the task is to use a balance scale to compare the total weights of any two subsets of items, aiming, after a fixed number of weighings, to partition the items into groups with as equal total weights as possible. \ouralgo improved the ALE-Agent's simulated annealing baseline by introducing faster caching, refining fallback weight estimation, and ultimately replacing simulated annealing with a more focused optimization combining greedy moves and targeted local search. Comparison with top human solutions suggests that for many tasks, there is ample room for improvement. Furthermore, often times \ouralgo tended to explore modifications staying close to the ALE-Agent's solution. This indicates the potential of overfitting to the initialization solution.